\documentclass[10pt,a4paper]{amsart}
\usepackage[margin=19mm]{geometry}
\usepackage{amsmath,amssymb,mathtools}
\usepackage[scr=rsfs,cal=euler]{mathalfa}
\usepackage{xfrac}
\usepackage[unicode,hidelinks,bookmarks=false]{hyperref}
\newcommand{\ie}{i.e.\ }
\newcommand{\cf}{cf.\ }
\newcommand{\resp}{respectively\ }
\newcommand{\Subsection}{\subsection}
\providecommand{\nobreakdash}{\nobreak\hbox{-}}
\setlength{\emergencystretch}{2em}
\multlinegap0pt
\usepackage{float}
\usepackage{mathtools}
\usepackage{amssymb}
\usepackage{bm}
\usepackage{xcolor}
\usepackage{algorithm}\usepackage{algpseudocode}
\newtheorem{lemma}{Lemma}
\newcommand{\LaP}{{\mathscr{L}}}
\newcommand{\mathd}{\mathrm{d}}
\title{An Efficient Laguerre Minimum Action Method for Computing Quasi-Potentials}
\author{Shenghe Huang, Yishuang Yue, Haijun Yu}
\date{Source: arXiv:2606.21114v1 (CC BY 4.0)}
\begin{document}
\maketitle

\noindent\emph{Source and licence.} An Efficient Laguerre Minimum Action Method for Computing Quasi-Potentials. Version arXiv:2606.21114v1 (CC BY 4.0), distributed by arXiv under the Creative Commons Attribution 4.0 International licence, http://creativecommons.org/licenses/by/4.0/, as recorded in the arXiv licence metadata for this exact version and shown on the abstract page https://arxiv.org/abs/2606.21114v1. Full licence notice: \texttt{licenses/arxiv-2606.21114-CC-BY-4.0.txt}.\par\smallskip

\noindent\textit{Editorial scope.} Counted unit: the article's lemma lemma (source0.tex lines 620--626). The article's complete original proof of it is delivered with it (source0.tex lines 628--654, one \texttt{proof} environment of 27 lines). No mathematics is written, repaired or re-derived: the statement and the proof are the article's own lines, in the article's own notation, and no retained line is substituted or re-flowed.  No further source-local context is needed: every cross-reference of every retained line resolves inside the two delivered spans.  Nothing was omitted from any retained span, so the package carries no omission record.  No retained line contains a citation, so the wrapper carries no bibliography and no bibliography entry.  No retained line contains a figure environment, an image-inclusion command or drawing code, so no image is delivered and the package declares no asset dependency.  Editorial formatting only, in addition: an AMS \texttt{amsart} preamble that restates the article's own declarations and macros used by the retained lines, namely the environments \texttt{lemma} on separate counters, together with the standard packages those lines need.  The retained environments print in this document as Lemma 1 in order of appearance, whereas the article prints the same items with the numbering of its own full document.  The complete native source of the article as distributed in the arXiv e-print for this exact version is bundled unchanged as \texttt{sources/203257/source0.tex}.
\begin{lemma}
	For $\varphi(t)=e^{-\widehat{Q}t}x_0$, assume that all eigenvalues of $\widehat{Q}$ have positive real parts. Then in the projection approximation, we have
	\[  \widehat{q }_n = \bigl( \frac{\hat{Q}}{\lambda} + \frac{1}{2} I
	\bigr)^{- 1} \bigl( I - \bigl( \frac{\hat{Q}}{\lambda} + \frac{1}{2} I
	\bigr)^{- 1} \bigr)^n x_0, \]
	where $I$ represents the identity matrix.
\end{lemma}

\begin{proof}
	Because of
	\[ \widehat{q }_n = \lambda \int^{\infty}_0 \varphi (t) \LaP_n (\lambda t)
	{e^{- \frac{\lambda t}{2}}}  \mathd t. \]
	Substitute $\varphi(t)=e^{-\widehat{Q}t}x_0$ into the above expression and $\lambda t$ into
	$\tau$, then we obtain
	\[ \widehat{q }_n = \int^{\infty}_0 \LaP_n (\tau) {e^{- \bigl(
			\frac{\hat{Q}}{\lambda} + \frac{1}{2} I \bigr) \tau}}  x_0 \mathd \tau .
	\]
	For the sake of simplicity, we denote $R = \frac{\hat{Q}}{\lambda} +
	\frac{1}{2} I$, and
	\[ \widehat{Q }_n = \int^{\infty}_0 \LaP_n (\tau) {e^{- R \tau}}  \mathd \tau,
	\]
	By the assumption, $R$ is invertible. The only thing we need to do is to
	calculate $\widehat{Q }_n$, and $\widehat{q }_n = \widehat{Q }_n x_0$. It is
	easy to get
	\[ \frac{\mathd^k}{\mathd \tau^k } \LaP_n (0) = (- 1)^k \binom{n}{n-k}, \quad k \in [0, n] \cap \mathbb{N} \]
	and
	\[ \frac{\mathd^{n + 1}}{\mathd \tau^{n + 1} } \LaP_n (0) = 0. \]
	By the integration by parts, we have
	\[ \int^{\infty}_0 \LaP_n (\tau) {e^{- R \tau}}  \mathd \tau = \sum_{k = 0}^n -
	R^{- (k + 1)} e^{- R t}  \frac{\mathd^k}{\mathd
		\tau^k } \LaP_n (\tau) \Big|_{\tau = 0}^{\infty}
	\]
	By the binomial theorem, we obtain
	\[ \widehat{Q}_k = R^{- 1} (I - R^{- 1})^k. \qedhere \] 
\end{proof}

\end{document}
